\subsection{Solutions des ateliers}
\begin{enumerate}
\item Les deux valeurs sont $3+2\ii$ et $-3-2\ii$ : leur carré vaut $9-4+12\ii=5+12\ii$. La factorisation de la différence des carrés exclut toute autre solution.
\item Le polynôme construit est $p(z)=z^{2} - z - 3 \mathrm{i} z - 4 + 3 \mathrm{i}$. Son discriminant est $8 - 6 \mathrm{i}=(3-\ii)^2$. En utilisant $\delta=3-\ii$, la formule quadratique redonne $2+\ii$ et $-1+2\ii$. Leur somme et leur produit contrôlent les deux coefficients.
\item La substitution donne $Z^2-(3-2\ii)Z+8+6\ii=(Z-2\ii)(Z-(3-4\ii))$. Comme $(1+\ii)^2=2\ii$ et $(2-\ii)^2=3-4\ii$, les solutions sont $\pm(1+\ii)$ et $\pm(2-\ii)$. Chaque facteur quadratique a exactement les deux racines indiquées ; les quatre nombres sont distincts et leur substitution annule le polynôme.
\end{enumerate}
